\documentclass[11pt,a4paper]{article}
\usepackage[T1]{fontenc}
\usepackage[utf8]{inputenc}
\usepackage[margin=22mm]{geometry}
\usepackage{lmodern}
\usepackage{amsmath,amssymb}
\usepackage{enumitem}
\usepackage{xcolor}
\usepackage{microtype}
\usepackage{hyperref}
\definecolor{epflred}{HTML}{E3000B}
\hypersetup{colorlinks=true,urlcolor=epflred,linkcolor=epflred,
  pdftitle={ENG-654 Project 4: Joint Limits and IK Distribution in Orthogonal 3R Robots}}
\setlength{\parindent}{0pt}
\setlength{\parskip}{6pt}
\setlist[itemize]{leftmargin=*,itemsep=5pt,topsep=4pt}
\urlstyle{tt}

\begin{document}
{\small\textbf{EPFL \textbar\ ENG-654}\hfill Kinematics-Grounded Motion Planning for Robots}
\vspace{5mm}

{\color{epflred}\Large\bfseries Project 4}

{\LARGE\bfseries Joint Limits and IK Distribution\\in Orthogonal 3R Robots\par}

\textbf{Format:} Group of two students; simulation-based project.\\
\textbf{Prerequisites:} D--H modelling, inverse kinematics, Jacobians,
singularities, joint limits, and configuration-space connectivity.

\section*{Motivation}
The mathematical workspace can overstate what a robot can execute. Joint
limits remove some postures and can disconnect a region that was connected
before those limits were imposed. Two target positions may each be reachable
yet inaccessible from the same initial configuration without violating a limit
or encountering a singularity.

\section*{Task}
\textbf{Overall objectives.} Compare how joint limits reshape the IK distribution and feasible
configuration components of selected cuspidal and noncuspidal orthogonal 3R
geometries.

\begin{itemize}
  \item \textbf{Establish two reference geometries.} Use the noncuspidal and cuspidal
idealized geometries: $a=(a_1,2,1.5)$ with $a_1=0$ and $1$ respectively,
$d=(1,1,0)$, and $\alpha=(\pi/2,\pi/2,0)$. Validate forward kinematics and the
position Jacobian, and keep these models distinct from the exact custom URDF.

  \item \textbf{Build the unconstrained maps.} Plot $\det J_p$ and its zero contours in
the $(q_2,q_3)$ plane. Map the singular curves into $(\rho,z)$ and identify
the two- and four-IK regions. Overlay the IK configurations of several targets
and record which sampled singularity-free component contains each solution.

  \item \textbf{Apply controlled limit boxes.} Compare the periodic unconstrained case
with $q_2,q_3\in[-150^\circ,150^\circ]$ and with
$q_2\in[-90^\circ,120^\circ]$, $q_3\in[-120^\circ,90^\circ]$.
Keep $q_1$ unrestricted so that the meridian reduction remains meaningful.
These are experimental limits, not manufacturer specifications.

  \item \textbf{Classify feasible inverse solutions.} Retain the original mathematical
IK count and separately count limit-admissible IKs, which may be one or three
as well as two or four. Recompute sampled configuration components after
clipping by limits. Do not connect opposite plot boundaries when that would
cross a physical joint limit.

  \item \textbf{Explain global accessibility.} For two diagnostic target pairs, compare
pointwise reachability with accessibility from a selected initial component.
Display component-labelled workspace coverage and validate one candidate
connection where available. Refine ambiguous bottlenecks and explain the
limitations of sampled connectivity claims.
\end{itemize}

\newpage
\section*{Specifics and scope}
Study two geometries and the three stated limit settings, giving six
comparable cases. Use the same workspace domain and color conventions for all
cases. Report a counterexample to equating pointwise reachability with path
accessibility if one is found; a carefully explained negative search is acceptable.
Collision modeling is outside the core experiment.

Use the position Jacobian $J_p$ for the 3R task.
Confirm that the meridian reduction is valid for the selected model and
operational point. Count distinct real IKs modulo periodicity, then filter
joint limits separately. Identify two-/four-IK regions where present and
explain any absent count rather than forcing both types into every map. Show both determinant values and the zero contour,
and distinguish sampled connectivity from a proof. Refine the decisive
region instead of claiming completeness from a coarse grid.

\section*{Expected outcome}
Distinguish mathematical IK multiplicity, limit-admissible solutions, and
accessibility within one feasible component. Deliver reduced determinant plots,
two-/four-IK reference maps, constrained IK-count maps, and a comparison of
component-dependent workspace coverage.

Submit reproducible code, model parameters and test data, the required plots,
a concise 5--6-page report, and a short simulation demonstration. Explain
both successful cases and failures; conclusions must follow the numerical
and geometric evidence rather than an expected outcome.

\section*{Provided materials and implementation}
The custom 3R robot description (.urdf) and link meshes (.stl) will be provided
to the students. The reference robot is orthogonal with non-intersecting
consecutive axes. It illustrates axis geometry; its dimensions must not be
substituted for those of the two idealized geometries.

Use standard D--H transforms
$T_{i-1,i}=R_z(q_i)T_z(d_i)T_x(a_i)R_x(\alpha_i)$, with zero joint-angle
offsets for the idealized models. Reuse a validated analytical 3R IK
implementation compatible with these parameters. Use plotting or robotics
tools to inspect frames, display IKs and sampled components, and animate
candidate connections; verify returned IKs with your forward model.

\end{document}
